
\documentclass[aspectratio=169]{beamer}

\mode<presentation>
\usetheme{Boadilla}
\definecolor{NTHU1}{RGB}{127,16,132}
\definecolor{NTHU2}{RGB}{228,210,231}
\definecolor{blue}{RGB}{30,90,205}
\definecolor{red}{RGB}{213,94,0}
\definecolor{green}{RGB}{0,128,0}
\definecolor{charcoalgray}{RGB}{108,104,104}
\setbeamercolor{title}{fg=NTHU1}
\setbeamercolor{frametitle}{fg=NTHU1}
\setbeamercolor{block title}{bg=NTHU1, fg=white}
\setbeamercolor{block body}{bg=NTHU2}
\setbeamercolor{structure}{fg=NTHU1}
\setbeamercolor{item projected}{fg=white}
\setbeamercolor{item}{fg=NTHU1}
\setbeamercolor{subitem}{fg=NTHU1}
\setbeamercolor{section in toc}{fg=NTHU1}
\setbeamercolor{description item}{fg=NTHU1}
\setbeamercolor{caption name}{fg=NTHU1}
\setbeamercolor{button}{bg=NTHU1, fg=white}
\setbeamercolor{caption name}{fg=NTHU1}
\usepackage{graphics}
\usepackage{tikz}
\usepackage{amsmath}
\usepackage{bbm}
\usetikzlibrary{decorations.pathreplacing}
\usepackage{geometry}
\usepackage{booktabs}
\usepackage{bookmark}
\usepackage{multirow, makecell}
\usepackage{float}
\usepackage{fancyvrb}
\usepackage{caption}
\captionsetup[figure]{singlelinecheck = false, format= hang, justification=centering}
\captionsetup[subfigure]{singlelinecheck = false, format= hang, justification=raggedright}
\usepackage{subcaption}
\usepackage{adjustbox}
\usepackage{hyperref}
\usepackage{threeparttable}
\usepackage{appendixnumberbeamer} 
\usepackage[T1]{fontenc}
\usepackage[default]{lato} 
\usepackage{smartdiagram}
\smartdiagramset{
  back arrow disabled = true,
  module minimum width=1cm,
  module x sep = 3,
  uniform arrow color=true,
}
\usepackage{xcolor}
\usepackage{colortbl}
  \definecolor{light-gray}{gray}{0.99}

\newenvironment{wideitemize}{\itemize\addtolength{\itemsep}{10pt}}{\enditemize}
\newenvironment{wideenumerate}{\enumerate\addtolength{\itemsep}{10pt}}{\endenumerate}
\newenvironment{widedescription}{\description\addtolength{\itemsep}{10pt}}{\enddescription}
\hypersetup{
colorlinks=true,
linkcolor=NTHU1,
filecolor=green, 
urlcolor=blue,
}
\beamertemplatenavigationsymbolsempty
\setbeamercolor{author in head/foot}{bg=white, fg=NTHU1}
\setbeamercolor{title in head/foot}{bg=white, fg=NTHU1}
\setbeamercolor{date in head/foot}{bg=white, fg=NTHU1}
\setbeamercolor{section in head/foot}{bg=white, fg=NTHU1}
\setbeamercolor{headline}{bg=white}
\setbeamertemplate{footline}{
    \leavevmode%
    \hbox{%
        \begin{beamercolorbox}[wd=.333333\paperwidth,ht=2.25ex,dp=1ex,center]{date in head/foot}%
            \usebeamerfont{date in head/foot}%\insertdate
        \end{beamercolorbox}%
        \begin{beamercolorbox}[wd=.333333\paperwidth,ht=2.25ex,dp=1ex,center]{title in head/foot}%
            \usebeamerfont{title in head/foot}\insertshorttitle
        \end{beamercolorbox}%
        \begin{beamercolorbox}[wd=.333333\paperwidth,ht=2.25ex,dp=1ex,right]{date in head/foot}%
            \usebeamerfont{date in head/foot}\insertframenumber{}\hspace*{2em}
        \end{beamercolorbox}}%
        \vskip0pt%
    }
    %% May need to script the introduction!!!!!!
%\setbeamercolor{page number in head/foot}{fg=NTHU1}
\setbeamertemplate{section in toc}[sections numbered]
\setbeamertemplate{subsection in toc}{\leavevmode\leftskip=3em\rlap{\hskip-1.75em\inserttocsectionnumber.\inserttocsubsectionnumber}
\inserttocsubsection\par}
\setbeamerfont{subsection in toc}{size=\footnotesize}
%\setbeamertemplate{headline}{%
  %\begin{beamercolorbox}[ht=5.5ex]{section in head/foot}
    %\vskip2pt\insertnavigation{0.33\paperwidth}\vskip2pt
  %\end{beamercolorbox}%
%}
\newenvironment{transitionframe}{\setbeamercolor{background canvas}{bg=white}\setbeamertemplate{footline}{} \begin{frame}[noframenumbering]}{\end{frame}}

\makeatletter
\let\@@magyar@captionfix\relax
\makeatother

\title[]{Conflict and Peace} % Change this regularly
\subtitle[]{Week 14: Personnel capacity of the state}
\author[Lee]{Seung-hun Lee }
\institute[]{Taipei School of Economics\\ National Tsing Hua University}

\date[]{December 5th, 2025}

\begin{document}
\begin{transitionframe}
\titlepage
\end{transitionframe}

%%% Color slides for section headers: Use for colloquium version (The ones bewteen \iffals and \fi)

\section{Overview of the topic}

\begin{frame}
\frametitle{Why does the public sector personnel matter?} 
Take up huge roles: In principle and in data
\begin{itemize}\setlength\itemsep{3pt}
\item As in other sectors, they embody human capital of the state
\item They determine provision of public services and implementation of policies
\item Takes up nearly 20\% of all employment  and 35\% of formal employment (World Bank)
\item Responsible for the public sector that takes up approximately 45\% of GDP (OECD)
\end{itemize}\medskip
However, the quality of public sector workers varies across countries
\begin{itemize}\setlength\itemsep{3pt}
\item 20-35\% of front-line workers (teachers, health care) absent in developing countries
\item Lack institutional structures to incentivize, monitor existing workers
\begin{itemize}\setlength\itemsep{3pt}
\item Poor structure $\to$ Absenteeism and inefficiencies $\to$ Bad public services (repeat)
\end{itemize}
\item Finding motivated and capable individuals to work is difficult
\begin{itemize}\setlength\itemsep{3pt}
\item Lower pay, more rigid mobility than private sector
\end{itemize}
\item Bureaucratic quality is also clustered with development 
\end{itemize}
\end{frame}

\begin{frame}
\frametitle{Public sector takes up huge share of employment globally} 
\begin{figure}
  \includegraphics[keepaspectratio, width=0.8\textwidth]{fig1.png}
\caption{Share of public sector employment in total employment, 2019}
\end{figure}
\begin{footnotesize}
Source: Our World in Data (\hyperlink{https://ourworldindata.org/grapher/public-sector-employment-as-a-share-of-total-employment?time=2019}{link}), reprocessed from World Bank and ILO data 
\end{footnotesize}
\end{frame}

\begin{frame}
\frametitle{Clustering of development and bureaucratic quality} 
\begin{figure}
  \includegraphics[keepaspectratio, width=0.6\textwidth]{fig2.png}
\end{figure}
\begin{footnotesize}
Source: Besley, Burgess, Khan, Xu (2022) ``Bureaucracy and Development", \textit{Annual Review of Economics}, 14, 397-424 
\end{footnotesize}
\end{frame}

\begin{frame}
\frametitle{Bureaucrat recruitment and principal-agent problem} 
Ideal bureaucrats vs reality
\begin{itemize}\setlength\itemsep{3pt}
\item High ability individuals with prosocial motives
\item But lower pay $\to$ smaller applicant pool, less qualified individuals
\item Mixed evidence from field experiments
\begin{itemize}
\item Motivated applicants (Dal Bó et al 2013) vs Applying `for money' (Desserano 2015)
\end{itemize}
\end{itemize}\medskip
Problem with incentivizing, monitoring existing public sector workers
\begin{itemize}\setlength\itemsep{3pt}
\item Public workers as agents, superiors (and citizens) as principal
\item In many cases, agents' interest not aligned with principals' interests
\begin{itemize}
\item Actions of agents are not always observable $\to$ pursue own interest
\end{itemize}
\item What structure algins interests while motivating agents is a big question
\end{itemize}\medskip
We look at various exploring questions on selection, retention, and motivation of personnel in public sector
\end{frame}

\begin{transitionframe}
\frametitle{Roadmap for today}
\tableofcontents
\end{transitionframe}

\section{Impact of leadership and politician capacities}
\subsection{Dal Bó, Finan, Folke, Persson, Rickne (2017): Who Becomes a Politician?}
\begin{transitionframe}
\begin{center}
  \begin{huge}
    \textcolor{NTHU1}{%
  Dal Bó, Finan, Folke, Persson, Rickne (2017):\\ Who Becomes a Politician?
    }
  \end{huge}
\end{center}
\end{transitionframe}

\begin{frame}
\frametitle{Q: How to attract competent leaders while being representative?} 
Can democracies form competent and representative leadership (inclusive meritocracy)?
\begin{itemize}\setlength\itemsep{3pt}
\item Voters want to elect leaders capable of efficiently implementing policies
\item Voters also want policy-makers who represent diverse interests
\item Classical economic models suggest that these two are in a trade-off relation
\begin{itemize}\setlength\itemsep{3pt}
\item Market wage < Returns from office for low-quality individuals (low opportunity cost)
\item Good politicians likely to `free-ride' on bad politicians $\to$ Supply of good politicians $\Downarrow$
\item Even if good candidates can be selected, may make broad representation difficult
\end{itemize}
\end{itemize}\medskip
This paper uses data on Swedish local politicians to identify patterns in inclusive democracy
\begin{itemize}\setlength\itemsep{3pt}
\item Includes elected and \textit{nonelected}: Address selection bias 
\item Detailed measures on quality and representation
\item Identifies how both competent and representative leaders are selected in democracies
\end{itemize}
\end{frame}

\begin{frame}
\frametitle{Context and data} 
Swedish National MPs and Municipal council members
\begin{itemize}\setlength\itemsep{3pt}
\item MPs are full-time, while most municipal council members have other work
\begin{itemize}
\item Mayors, selected among council members, are full-time and paid highly
\end{itemize}
  \item Low monetary cost to individuals running for local office
  \item Key motives: Intrinsic concerns with policy, social interaction in policy circles
  \item Left represents blue-collar, center-right represents white-collar workers
\end{itemize}\medskip
Data obtained from 
\begin{itemize}\setlength\itemsep{3pt}
\item  Sample: All candidates from 1980-2010 
\item Links candidates to admin registers (including military enlistment and tax records)
\item Age, sex, education level, earnings, occupation, parental/sibling info, cognitive capacity 
\end{itemize}
\end{frame}

\begin{frame}
\frametitle{Positive selection of politicians based on abilities} 
\begin{figure}
  \includegraphics[keepaspectratio, width=0.65\textwidth]{dffpr2017_f1.png}
\end{figure}
\begin{itemize}
\item Beats classical prediction that there are more low-quality individuals entering politics
\end{itemize}
\end{frame}


\begin{frame}
\frametitle{Are they socially representative?}
Determining source of positive selection
\begin{itemize}\setlength\itemsep{3pt}
\item Privileged access? Ability should matter little once elite status is controlled for %$ By the heirs of rich families
\item Exclusive meritocracy? Key abilities are strongly associated with social background
\item Inclusive meritocracy? High ability individuals come from all social backgrounds
\end{itemize}\medskip
Empirical tests
\begin{itemize}\setlength\itemsep{3pt}
\item Compare with siblings: Same "elitist background" $\to$ Positive selection still appears
\item Parental income: Social background of politicians is evenly distributed
\item Social class: Politicians from wide range of occupations represented
\end{itemize}
\end{frame}

\begin{frame}
\frametitle{Elitism unlikely to explain positive selection} 
\begin{figure}
  \includegraphics[keepaspectratio, width=0.7\textwidth]{dffpr2017_f2.png}
\end{figure}
\end{frame}

\begin{frame}
\frametitle{Neither does exclusive meritocracy} 
\begin{figure}
  \includegraphics[keepaspectratio, width=0.8\textwidth]{dffpr2017_f3.png}
\end{figure}
\end{frame}

\begin{frame}
\frametitle{Discussion and takeaways}
What drives competence and representation? 
\begin{itemize}\setlength\itemsep{3pt}
\item Weak trade-off between competence and social background
\item Positive selection on ability even among those with poorer backgrounds
\item Intrinsic motives: Positive selection even for low-paid municipal councils
\item Screening by parties: Sorts more able people to higher positions
\end{itemize}\medskip
Remaining questions
\begin{itemize}\setlength\itemsep{3pt}
\item When and where do these findings hold outside 1980-2010 Sweden?
\begin{itemize}\setlength\itemsep{3pt}
\item Intrinsically motived individuals may not enter politics elsewhere, how to attract them?
\item Party incentives: Promote someone who can win vs. who has best policy
\item Other democratic countries? Sweden now? 
\end{itemize}
\end{itemize}
\end{frame}

\subsection{Xu (2018): The Costs of Patronage: Evidence from the British Empire}
\begin{transitionframe}
\begin{center}
  \begin{huge}
    \textcolor{NTHU1}{%
  Xu (2018):\\ The Costs of Patronage: Evidence from the British Empire
    }
  \end{huge}
\end{center}
\end{transitionframe}

\begin{frame}
\frametitle{Q: How does patronage affect allocation of positions \& performance?} 
How personnel are promoted and incentivized determines organizational performance
\begin{itemize}\setlength\itemsep{3pt}
\item Past: patronage widely used to appoint public sector positions
\begin{itemize}
\item Discretionary appointments based on network and ties to the ruler
\end{itemize}
\item In theory, discretion can align incentives between principal and angents
\item However, it can worsen performance if favoritism disincentivizes capable subordinates
\end{itemize}\medskip
This paper uses historical data to identify how patronage affects performance
\begin{itemize}\setlength\itemsep{3pt}
\item Colonial governors of British empire: 1854-1966 across 70 colonies
\item Connection determined by genealogy and biological data
\item Within-governor variation: Change in secretary of state $\to$ Change in connections
\item Across-time variation: Discretion (1854-1930) vs. Independent appointment (1930 -)
\item Test who gets assigned to richer colonies and how tax returns are affected
\end{itemize}

\end{frame}

\begin{frame}
\frametitle{Organization of British colonies}
Assignment of colonial governors 
\begin{itemize}\setlength\itemsep{3pt}
\item Handled by Colonial Office under Secretary of State for the Colonies in London
\item Secretary of State in charge for around 3 years
\item 1930 Reform replace Secretary of State's discretion with open recruitment
\end{itemize}\medskip
Role of the governors
\begin{itemize}\setlength\itemsep{3pt}
\item Fixed 6-year term
\item Responsible for public finance, legislation, and personnel appointments
\item Discretion over executive and finances: Widely different policies across colonies
\item Direct monitoring impossible due to distance
\end{itemize}\medskip
\end{frame}

\begin{frame}
\frametitle{Data and empirical strategy}
Digitalized (by Guo Xu himself!) archival data 
\begin{itemize}\setlength\itemsep{3pt}
\item Sources: Colonial Blue Books, Colonial Office List, Who's Who, The Peerage
\item Postings (and their wages), background of governors, genealogy 
\begin{itemize}
\item Connection determined by ancestry, membership at aristocratic clubs, and schools
\end{itemize}
\item Colonial statistics (revenue, expenditure, demographics etc)
\end{itemize}\medskip
Empirical strategy leverages differences in connectedness across different periods\vspace{-10pt}
\[
y_{it} = \beta c_{it}+\theta_i + \gamma X_{it}+\tau_t + e_{it}
\]\vspace{-15pt}
\[
y_{it} = \beta_0 c_{it} + \beta_1 c_{it}\times I[t\geq1930]+\theta_i + \gamma X_{it}+\tau_t + e_{it}
\]\vspace{-10pt}
\begin{itemize}\setlength\itemsep{3pt}
\item Connection $c_{it}$ used with various criterion
\item $c_{it}$ is predetermined and driven by secretary of state changes
\begin{itemize}
  \item No reverse causality: Unlikely to attend Oxbridge to be connected to current governor
  \item Secretary of state turnover generates common shocks to all serving governors
\end{itemize}
\end{itemize}
\end{frame}

\begin{frame}
\frametitle{Connected governor salaries are higher! } 
    \begin{figure}
    \includegraphics[keepaspectratio, width=0.8\textwidth]{x2018_t1.png}
    \end{figure}
\end{frame}

\begin{frame}
\frametitle{Salary gap washed away post-1930} 
    \begin{figure}
    \includegraphics[keepaspectratio, width=0.7\textwidth]{x2018_f1.png}
    \end{figure}
\end{frame}

\begin{frame}
\frametitle{Patronage worsens performance, which is offset following reforms} 
\begin{columns}
  \begin{column}{0.65\textwidth}
    \begin{figure}
    \includegraphics[keepaspectratio, width=\textwidth]{x2018_t2_1.png}
    \end{figure}
  \end{column}
  \begin{column}{0.3\textwidth}
    \begin{figure}
       \includegraphics[keepaspectratio, width=\textwidth]{x2018_t2_2.png}
       \end{figure}
  \end{column}
\end{columns}
\end{frame}

\begin{frame}
\frametitle{Takeaways and discussions} 
What the findings say about organizational performance
\begin{itemize}\setlength\itemsep{3pt}
\item Connection $\to$ Positions that yield higher wages
\item Lower revenue collection: Due to higher tax exemptions, high reliance on customs
\item Reforms washed salary gaps and shortfalls in performance
\end{itemize}\medskip
Discussions
\begin{itemize}\setlength\itemsep{3pt}
\item Favoritism can lead to worse organizational performance, even in public institutions
\item But how does it affect recruitment into the pool of governor candidates?
\begin{itemize}\setlength\itemsep{3pt}
\item If people know connections matter, talented individuals with no connection discouraged
\item Leaving pool of candidates with less talented individuals
\item Ways to look for `outsiders' with talent?
\end{itemize}
\end{itemize}
\end{frame}

\subsection{Fenizia (2022): Managers and Productivity in the Public Sector}
\begin{transitionframe}
\begin{center}
  \begin{huge}
    \textcolor{NTHU1}{%
  Fenizia (2022):\\ Managers and Productivity in the Public Sector
    }
  \end{huge}
\end{center}
\end{transitionframe}

\begin{frame}
\frametitle{Q: Do talented managers in public sector improve performance?} 
How important are managers in public sectors?
\begin{itemize}\setlength\itemsep{3pt}
\item Oversee day-to-day operations and supervise policy implementation 
\item Limited tools available compared to private sector managers`'
\begin{itemize}
\item Limited discretion in firing and promoting workers
\end{itemize}
\item How to measure performance of government agencies?
\end{itemize}\medskip
This paper uses innovative measures of performance to study the role of managers
\begin{itemize}\setlength\itemsep{3pt}
\item Italian Social Security Agency data (INPS): Covering varous welfare programs
\item Uses welfare claims as measure of output (hoogenous across different offices)
\item All sites have same rules and physical capital$\to$ distinguish manager contributions
\end{itemize}\medskip
\end{frame}

\begin{frame}
\frametitle{Background and data} 
Role of INPS 
\begin{itemize}\setlength\itemsep{3pt}
\item Administers all social welfare and insurance programs in Italy 
\item Claims are processed in local offices run by single managers 
\item Managers allocate work, determine workplace practices, but cannot fire workers
\item Managers selected among limited pool of candidates and rotate every 5 years 
\end{itemize}\medskip
Data sources on managers (2011Q1 - 2017Q2)
\begin{itemize}\setlength\itemsep{3pt}
\item Output: $\sum$ claims weighed by complexity (procesing time) per worker
\item Service quality: Timeliness and error rotate
\item Worker characteristics: Job title, age, office locations etc
\end{itemize}
\end{frame}

\begin{frame}
\frametitle{Empirical strategies} 
Do productive managers improve performance?
\begin{itemize}\setlength\itemsep{3pt}
\item Fixed effects regression with office($i$), time($t$), and manager($m(i,t)$) fixed effects
\[
\log{P}_{it} = \alpha_i + \tau_t + \theta_{m(i,t)}+u_{it}
\]\vspace{-20pt}
\item Requires that manager moves are exogenous conditional on office and time effects
\begin{itemize}
  \item Observable characteristics should not predict manager quality
\item Violated if better managers are assigned to worse offices
\end{itemize}
\item Compare adjusted $R^2$ to test how much of output variation is explained by managers 
\end{itemize}\medskip
What makes for a productive manager?
\begin{itemize}\setlength\itemsep{3pt}
\item Regress changes in outcomes with quarter before manager change to talent changes
\[
y_{i,t+k}-y_{i,t-1} = \pi_0^k + \pi_1^k(\hat{\theta}_{i,new}^k-\hat{\theta}_{i,old})+\gamma^k X_{i}+ (e_{i,t+k}-e_{i,t-1})
\]\vspace{-20pt}

\end{itemize}
\end{frame}

\begin{frame}
\frametitle{Managers explain 10\% of variation in productivity } 
\begin{figure}
  \includegraphics[keepaspectratio, width=0.57\textwidth]{f2022_t1.png}
\end{figure}
\begin{itemize}
\item Adjusted $R^2$: 0.68 in (2) $\to$ 0.76 in (3): 8 p.p increase
\item $0.08/0.762\simeq 0.104$: Manager explains 10\% of variation in productivity
\item Robust to selective sorting, other ways of variance decomposition
\end{itemize}

\end{frame}


\begin{frame}
\frametitle{Talented managers improve productivity, with workers decreasing...} 
\begin{figure}
  \includegraphics[keepaspectratio,width=0.6\textwidth]{f2022_f1.png}
\end{figure}
\end{frame}

\begin{frame}
\frametitle{Talented managers induced by retirement of older workers} 
\begin{figure}
  \includegraphics[keepaspectratio,width=0.6\textwidth]{f2022_t2.png}
\end{figure}
\end{frame}

\begin{frame}
\frametitle{Takeaways and discussions}
Other outcomes 
\begin{itemize}\setlength\itemsep{3pt}
\item Trade-off in quality? Not observed
\item Outcomes not driven by shifting towards easier claims to handle
\item Thus managers do have a meaningful impact on productivity
\end{itemize}\medskip
Remaining questions
\begin{itemize}\setlength\itemsep{3pt}
\item The managers oversee relatively small teams with well-defined goals this context?
\begin{itemize}\setlength\itemsep{3pt}
\item What about large projects and teams?
\item What about organization with no fixed roles and less structure?
\end{itemize}
\item Other factors may determine managers' role in productivity
\begin{itemize}\setlength\itemsep{3pt}
\item Ability to obtain inputs efficiently (no bribing, better communication)
\item Ability to attract more talented workers under the command 
\end{itemize}
\end{itemize}\medskip
\end{frame}


\subsection{Jones, Olken (2005): Do Leaders Matter? National Leadership and Growth since World War II [student presentation]}
\begin{transitionframe}
\begin{center}
  \begin{huge}
    \textcolor{NTHU1}{%
  Jones, Olken (2005): \\ Do Leaders Matter? National Leadership and Growth since World War II [student presentation]
    }
  \end{huge}
\end{center}
\end{transitionframe}




\section{Recruiting, incentivizing, and monitoring bureaucrats}
\subsection{Dal Bó, Finan, Rossi (2013): Strengthening State Capabilities: The Role of Financial Incentives in the Call to Public Service}
\begin{transitionframe}
\begin{center}
  \begin{huge}
    \textcolor{NTHU1}{%
  Dal Bó, Finan, Rossi (2013):\\ Strengthening State Capabilities: The Role of Financial Incentives in the Call to Public Service
    }
  \end{huge}
\end{center}
\end{transitionframe}

\begin{frame}
\frametitle{Q: Can fiscal incentives attract capable individuals to public sector?} 
Professionalized bureaucracy require individuals with capabilities and integrity
\begin{itemize}\setlength\itemsep{3pt}
\item These do not always co-move
\begin{itemize}\setlength\itemsep{3pt}
\item Capable individuals may seek opportunities in private sectror (higher pay)
\item This may leave public sector with less able individuals
\item Not all capable individuals have integrity:
\end{itemize}
\item So what makes up motivated agents and how to attract them to public sector?
\end{itemize}\medskip
This paper investigates role of incentives while investigating various aspects of quality
\begin{itemize}\setlength\itemsep{3pt}
\item How do monetary and non-monetary incentives attract candidates?
\item Estimates if higher wages attract more candidates and those of higher quality
\item Also estimates the labor supply elasticity facing public sector in Mexico
\item Heterogeneity by locational characteristics (local crime rates)
\end{itemize}
\end{frame}



\begin{frame}
\frametitle{Context of the experiment}
Mexico's regional development program (RDP) in 2011
\begin{itemize}\setlength\itemsep{3pt}
  \item Network of 350 agents \& 50 coordinators to identify needs, report to federal govt
\item Exogenously assign different wage offers across different posts (5000 vs 3750 MXP)
\item Candidates screerned to document various characteristics
\begin{itemize}\setlength\itemsep{3pt}
  \item Aptitude: intelligence  (IQ) and other qualities known to determine market value
  \item Motivation: Integrity, prosocial inclinations, public service motivation (interviews)
  \item Attraction \& commitment to policy making, justice, civicness, compassion, self-sacrifice
\end{itemize}
\item Vacancies randomized to one of 4 strata $\to$ Offers randomized to indiv. in that stratum
\begin{itemize}\setlength\itemsep{3pt}
  \item High vs low wage, High ($>9$) vs normal IQ (7-9 Ravens matrix)
\end{itemize}
\end{itemize}\medskip
Estimating causal effects: OLS (due to random assignment design)
\begin{itemize}\setlength\itemsep{3pt}
  \item Effect on character $Y$ in applicant pool (region): $Y_{icr}=\beta_1 T_c + \xi_r + e_{icr}$
  \item Effect on acceptance $A$ among applicants (strata): $A_{ics}=\gamma_1 T_c + \beta X_i + \xi_s + e_{ics}$
  \item Hetergeneity by municipalities: $A_{ics}=\gamma_1 T_c + \gamma_2 T_c \times W_m + \gamma_3 W_m + \beta X_i + \xi_s + e_{ics}$
\end{itemize}
\end{frame}

\begin{frame}
\frametitle{Higer wage $\to$ High ability applicants $\Uparrow$} 
\begin{columns}
\begin{column}{0.5\textwidth}
\begin{figure}
  \includegraphics[keepaspectratio, width=\textwidth]{dfr2013_t1_1.png}
\end{figure}
\end{column}
\begin{column}{0.5\textwidth}
\begin{figure}
  \includegraphics[keepaspectratio, width=\textwidth]{dfr2013_t1_2.png}
\end{figure}
\end{column}
\end{columns}
\end{frame}

\begin{frame}
\frametitle{Higer wage $\to$ Better motivated applicants $\Uparrow$} 
\begin{columns}
\begin{column}{0.5\textwidth}
\begin{figure}
  \includegraphics[keepaspectratio, width=\textwidth]{dfr2013_t2_1.png}
\end{figure}
\end{column}
\begin{column}{0.5\textwidth}
\begin{figure}
  \includegraphics[keepaspectratio, width=\textwidth]{dfr2013_t2_2.png}
\end{figure}
\end{column}
\end{columns}
\end{frame}

\begin{frame}
\frametitle{Higer wage $\to$ Applicants more likely to accept} 
\begin{figure}
  \includegraphics[keepaspectratio, width=0.7\textwidth]{dfr2013_t3.png}
\end{figure}
\end{frame}

\begin{frame}
\frametitle{Higher wage needed to offset disadvantageous regional traits} 
\begin{figure}
  \includegraphics[keepaspectratio, width=0.7\textwidth]{dfr2013_t4.png}
\end{figure}
\end{frame}

\begin{frame}
\frametitle{Additional findings and takeaways}
Labor supply elasticity $(\epsilon)$ for public sector with respect to wages 
\begin{itemize}\setlength\itemsep{3pt}
\item $\epsilon  = \epsilon_{\text{pool}}  + \epsilon_{\text{accept}}+\epsilon_{\text{pool}} \times \epsilon_{\text{accept}}\times\text{wage difference}$
\item Wage difference: 33\% (=(5000-3750)/3750)
\item Pool: 4.7 person $\uparrow$ from 18.1 person (26.3\%$\uparrow$) (Table 3) $\to$ Elasticity: 0.8 (=26.3/33)
\item Acceptance:  15.1 p.p $\uparrow$ from 42.9\% (35.2\% $\uparrow$) (Table 5) $\to$ Elasticity: 1.07 (=35.2/33)
\item $0.8+1.07+0.8\times1.07\times0.33=2.15$ ( 1\% $\uparrow$ in wages $\to$ Supply  2.15\% $\uparrow$)
\end{itemize}\medskip
Takeaways
\begin{itemize}\setlength\itemsep{3pt}
\item Offering high wages attracts better and motivated candidates
\item However, findings hold only if motivation and ability are positively correlated
 \begin{itemize}
\item What if service commitment is sacrificed for higher pay (prioritize superiors over citizens?)
\end{itemize}
\item Silent on post-acceptance performance: High pay to keep people motivated once hired? 
\end{itemize}
\end{frame}

\subsection{Bandiera, Best, Khan, Prat (2021): The Allocation of Authority in Organizations: A Field Experiment with Bureaucrats}
\begin{transitionframe}
\begin{center}
  \begin{huge}
    \textcolor{NTHU1}{%
  Bandiera, Best, Khan, Prat (2021):\\ The Allocation of Authority in Organizations: A Field Experiment with Bureaucrats
    }
  \end{huge}
\end{center}
\end{transitionframe}

\begin{frame}
\frametitle{Q: How does allocation of authority to bureaucrats affect outcomes?} 
Organization bring together people with different interests and skills
\begin{itemize}\setlength\itemsep{3pt}
\item Working towards a common goal challenging: Principal-agent problem and monitoring
\item How to allocate decision rights at different layers of hierarchy?
\item How to monitor and incentivize agents' behavior?
\item Many work on effects of performance pay, few look at decision-making structure
\end{itemize}\medskip
This paper designs an experiment in Pakistan procurement sector that combines the two
\begin{itemize}\setlength\itemsep{3pt}
\item Strict rules and monitoring vs. simplification and more autonomy?
\item Procurement officers (PO, make purchases) and Monitors (approve purchase) 
\begin{itemize}\setlength\itemsep{3pt}
\item Monitors: Approve purchases quickly vs Delay purchase until end of fiscal year
\end{itemize}
\item Experiment: Control + Autonomy + Performance pay + Both 
\item Different policy instrument works in different structures
\end{itemize}
\end{frame}

\begin{frame}
\frametitle{Procurement structure in Pakistan and theoretical framework} 
Procurement structure
\begin{itemize}\setlength\itemsep{3pt}
\item Experiment took place in 26 / 36 districts in Punjab (80\% of population)
\item Each office has 1 PO with allocated budgets
\item POs submit purchase for pre-audit approval to a monitor in each district
\begin{itemize}\setlength\itemsep{3pt}
\item This monitor can ask for extra paperwork (withhold) or approve the purchase
\end{itemize}
\end{itemize}\medskip
Conceptual framework: Expected effects of autonomy vs incentives by type of monitors
\begin{itemize}\setlength\itemsep{3pt}
\item Aligned PO and monitors behave in the interest of citizens (low price, little delays)
\item Misaligned PO raises price (bribes, little effort to find low-price)
\item Misaligned monitor delays purchase (exert power on PO)
\item Autonomy: Price$\downarrow$ if PO more aligned then monitor (opposite if PO more misaligned)
\item Incentives: Motivate PO, but dampened with misaligned monitors (higher costs)
\end{itemize}
\end{frame}

\begin{frame}
\frametitle{Experimental design} 
Measuring bureaucratic performance with value for money (with online platform)
\begin{itemize}\setlength\itemsep{3pt}
\item Price conditional on quantity, nature of purchased goods, and net of transaction costs
\begin{itemize}\setlength\itemsep{3pt}
\item Items: Consumables such as papers, toners, pen, banners, light bulbs etc
\item Price is adjusted for delivery speed, cost of delivery
\end{itemize}
\item Others: Expected price, good-bad purchase categorization(based on attributes), ML
\end{itemize}\medskip
Experimental design
\begin{itemize}\setlength\itemsep{3pt}
\item Autonomy: Remove the role of monitors in purchasing decisions
\item Incentive: Range of bonus for improving value for money ($\times2$  - $\times1/2$ of monthly wage)
\item Another group that receives both
\item Stratified by district-department level\vspace{-10pt}
 \[
 y_{igto} = \alpha + \sum_{k=1}^3\eta_k \text{Treatment}_o^k + \rho_g q_{igto}+\delta_s +\gamma_g + e_{igto}
 \]
\end{itemize}
\end{frame}

\begin{frame}
\frametitle{Overall, autonomy $\to$ lower prices paid but not incentives} 
\begin{figure}
  \includegraphics[keepaspectratio, width=\textwidth]{bbkp2021_t1.png}
\end{figure}
\end{frame}

\begin{frame}
\frametitle{Effects differ by type of monitors} 
\begin{figure}
  \includegraphics[keepaspectratio, width=0.7\textwidth]{bbkp2021_f1.png}
\end{figure}
\begin{itemize}\setlength\itemsep{3pt}
\item Better monitors (low delays, AG June share $\downarrow$): Incentive lower price (not autonomy)
\item Worse monitors (more delays, AG June share $\uparrow$): Autonomy lower price (not incentives)
\end{itemize}
\end{frame}

\begin{frame}
\frametitle{Treatment $\to$ large savings, but heterogenous by monitor types} 
\begin{figure}
  \includegraphics[keepaspectratio, width=0.7\textwidth]{bbkp2021_f2.png}
\end{figure}
\end{frame}

\begin{frame}
\frametitle{Takeaways and remaining questions} 
Incentives and autonomy work in different situations
\begin{itemize}\setlength\itemsep{3pt}
\item Incentives to low-level agents mean little if their supervisors are misaligned
\item In such case, giving low-level agents more autonomy improves performance
\item If supervisors are aligned, incentives improve effort of angents
\item On the other hand, autonomy does not change much from status quo in this case
\end{itemize}\medskip
Remaining questions on the organizational Economics
\begin{itemize}\setlength\itemsep{3pt}
\item Can stringent rules, not autonomy, be an optimal choice?
\begin{itemize}\setlength\itemsep{3pt}
\item Could still be worthwhile to prevent misaligned agents imposing high costs
\end{itemize}
\item How does autonomy affect \textit{who becomes a bureaucrat}?
\begin{itemize}\setlength\itemsep{3pt}
\item How to bring in more aligned agents (make efficient decisions for public) than misaligned agents (extract personal benefits with autonomy)?
\end{itemize}
\end{itemize}\medskip
\end{frame}

\subsection{Callen, Gulzar, Hasanain, Khan, Rezaee (2023): The Political Economy of Public Sector Absence [student presentation]}
\begin{transitionframe}
\begin{center}
  \begin{huge}
    \textcolor{NTHU1}{%
  Callen, Gulzar, Hasanain, Khan, Rezaee (2023):\\ The Political Economy of Public Sector Absence [student presentation]
    }
  \end{huge}
\end{center}
\end{transitionframe}

\begin{frame}
\frametitle{Takeaways and remaining questions} 
Takeaways
\begin{itemize}\setlength\itemsep{3pt}
\item Better leaders and bureaucrats improve performance of organizations
\item Personnel selected on factors other than merit can hurt performance
\item Different organizational structure works best to align principal-agent interests
\begin{itemize}\setlength\itemsep{3pt}
\item If lower-ranked agents are more aligned than middle-managers, more autonomy helps
\item If different levels of bureaucracy are similarly aligned, incentives induces better work
\end{itemize}
\item Recruiting capable \& motivated agents, while monitoring \& incentivizing workers are difficult to balance
\end{itemize}\medskip
Remaining questions
\begin{itemize}\setlength\itemsep{3pt}
\item Delineation of roles left to `politicians' and `bureaucrats'
\begin{itemize}\setlength\itemsep{3pt}
\item Perform different roles and subject to different accountability measures
\item What roles are better left to politicians vs bureaucrats?
\end{itemize}
\item Relationship with private sector
\begin{itemize}\setlength\itemsep{3pt}
\item Complement developmental efforts vs limit/capture each other?
\end{itemize}
\end{itemize}
\end{frame}


%%%%%%%%%%%
\end{document}


